\newpage
\phantomsection
\label{prf:continuity-algebra-at-a-point}

\begin{remark*}[Return]
\hyperref[thm:continuity-algebra-at-a-point]{Return to Theorem}
\end{remark*}

\begin{theorem*}[Algebra of Continuity at a Point]
Let $A \subseteq \mathbb{R}$, let $f,g : A \to \mathbb{R}$, let $b \in \mathbb{R}$, and let $c \in A$. If $f$ and $g$ are continuous at $c$, then $f+g$, $f-g$, $fg$, $bf$ are continuous at $c$, and $f/g$ is continuous at $c$ whenever $g(c) \neq 0$.
\end{theorem*}

\begin{proof}
\textbf{Professional Standard Proof.}~\\
Let $A \subseteq \mathbb{R}$, let $f,g : A \to \mathbb{R}$, let $b \in \mathbb{R}$, and let $c \in A$. Assume that $f$ and $g$ are continuous at $c$. Since $f$ and $g$ are continuous at $c$, one has $\lim_{x\to c} f(x)=f(c)$ and $\lim_{x\to c} g(x)=g(c)$. By the sum rule for limits, $\lim_{x\to c} (f+g)(x)=f(c)+g(c)=(f+g)(c)$, so $f+g$ is continuous at $c$. By the difference rule, $\lim_{x\to c} (f-g)(x)=f(c)-g(c)=(f-g)(c)$, so $f-g$ is continuous at $c$. By the product rule, $\lim_{x\to c} (fg)(x)=f(c)g(c)=(fg)(c)$, so $fg$ is continuous at $c$. By the scalar multiple rule, $\lim_{x\to c} (bf)(x)=b\,f(c)=(bf)(c)$, so $bf$ is continuous at $c$. For the quotient, assume $g(c) \neq 0$. Then $\lim_{x\to c} g(x)=g(c) \neq 0$, so by the quotient rule for limits, $\lim_{x\to c} \frac{f(x)}{g(x)} = \frac{f(c)}{g(c)} = \left(\frac{f}{g}\right)(c)$, hence $f/g$ is continuous at $c$.
\end{proof}

\begin{proof}
\textbf{Detailed Learning Proof.}~\\

\textbf{Step 1.} Translate continuity hypotheses into limit form. Since $f$ and $g$ are continuous at $c$, the definition of continuity at a point gives $\lim_{x\to c} f(x)=f(c)$ and $\lim_{x\to c} g(x)=g(c)$. This reformulation allows application of limit algebra rules.

\textbf{Step 2.} Establish continuity of $f+g$ at $c$. By the sum rule for limits of functions, $\lim_{x\to c} (f+g)(x) = \lim_{x\to c} f(x) + \lim_{x\to c} g(x) = f(c) + g(c) = (f+g)(c)$. This limit equation is precisely the definition of continuity of $f+g$ at $c$.

\textbf{Step 3.} Establish continuity of $f-g$ at $c$. By the difference rule for limits of functions, $\lim_{x\to c} (f-g)(x) = \lim_{x\to c} f(x) - \lim_{x\to c} g(x) = f(c) - g(c) = (f-g)(c)$. Hence $f-g$ is continuous at $c$.

\textbf{Step 4.} Establish continuity of $fg$ at $c$. By the product rule for limits of functions, $\lim_{x\to c} (fg)(x) = \left(\lim_{x\to c} f(x)\right) \left(\lim_{x\to c} g(x)\right) = f(c)g(c) = (fg)(c)$. Hence $fg$ is continuous at $c$.

\textbf{Step 5.} Establish continuity of $bf$ at $c$. By the scalar multiple rule for limits of functions, $\lim_{x\to c} (bf)(x) = b \lim_{x\to c} f(x) = b \cdot f(c) = (bf)(c)$. Hence $bf$ is continuous at $c$.

\textbf{Step 6.} Establish continuity of $f/g$ at $c$ when $g(c) \neq 0$. The hypothesis $g(c) \neq 0$ combined with $\lim_{x\to c} g(x) = g(c)$ ensures that the denominator remains nonzero in a neighborhood of $c$. By the quotient rule for limits of functions, $\lim_{x\to c} \frac{f(x)}{g(x)} = \frac{\lim_{x\to c} f(x)}{\lim_{x\to c} g(x)} = \frac{f(c)}{g(c)} = \left(\frac{f}{g}\right)(c)$. Hence $f/g$ is continuous at $c$.
\end{proof}

\begin{remark*}[Proof structure]
This proof employs a direct approach, translating the continuity hypotheses into limit form and then systematically applying the algebra rules for limits of functions. The strategy works because continuity at a point is equivalent to the limit equation $\lim_{x \to c} f(x) = f(c)$, which transforms the problem into verifying limit equations for the combined functions using established limit algebra.
\end{remark*}

\begin{remark*}[Dependencies]~\\
\begin{itemize}
  \item \hyperref[def:continuous-at-a-point]{Definition of continuity at a point}
  \item \hyperref[thm:limit-function-sum]{Sum Rule for Limits of Functions}
  \item \hyperref[thm:limit-function-difference]{Difference Rule for Limits of Functions}
  \item \hyperref[thm:limit-function-product]{Product Rule for Limits of Functions}
  \item \hyperref[thm:limit-function-quotient]{Quotient Rule for Limits of Functions}
  \item \hyperref[thm:limit-function-scalar]{Scalar Multiple Rule for Limits of Functions}
\end{itemize}
\end{remark*}
